\documentclass[11pt,a4paper]{article}

% ---------- Encodage et langue ----------
\usepackage[utf8]{inputenc}
\usepackage[T1]{fontenc}
% \usepackage{lmodern}  % décommenter si disponible (police vectorielle)
\usepackage[french]{babel}
\usepackage[a4paper,margin=2.3cm]{geometry}

% ---------- Mathématiques ----------
\usepackage{amsmath,amssymb,amsthm}
\usepackage{mathtools}
\usepackage{enumitem}
\usepackage[most]{tcolorbox}
\usepackage[colorlinks=true,linkcolor=blue!60!black,urlcolor=blue!60!black]{hyperref}

% ---------- Macros ----------
\newcommand{\R}{\mathbb{R}}
\newcommand{\N}{\mathbb{N}}
\newcommand{\Z}{\mathbb{Z}}
\newcommand{\Q}{\mathbb{Q}}
\newcommand{\C}{\mathbb{C}}
\newcommand{\K}{\mathbb{K}}
\newcommand{\Sun}{\mathbb{S}^1}
\newcommand{\interieur}[1]{\mathring{#1}}
\newcommand{\adh}[1]{\overline{#1}}
\newcommand{\Int}{\operatorname{int}}
\newcommand{\diam}{\operatorname{diam}}
\newcommand{\dist}{\operatorname{dist}}
\newcommand{\Vect}{\operatorname{Vect}}
\newcommand{\Bf}{B_f}
\newcommand{\norm}[1]{\left\|#1\right\|}
\newcommand{\abs}[1]{\left|#1\right|}
\newcommand{\opnorm}[1]{|\!|\!|#1|\!|\!|}
\newcommand{\ps}[2]{\left\langle #1,#2\right\rangle}
\newcommand{\Lc}{\mathcal{L}}
\newcommand{\Cz}{\mathcal{C}^0}
\newcommand{\Cinf}{\mathcal{C}^\infty}
\newcommand{\Cc}{\mathcal{C}^0_c}
\newcommand{\ie}{c'est-à-dire}
\DeclareMathOperator{\Ker}{Ker}
\DeclareMathOperator{\Ima}{Im}
\DeclareMathOperator{\supp}{supp}
\DeclareMathOperator{\sgn}{sgn}
\renewcommand{\Re}{\operatorname{Re}}

% ---------- Environnements ----------
\newtcolorbox{exo}[1]{
  enhanced, breakable,
  colback=black!3, colframe=black!70,
  fonttitle=\bfseries, title={#1},
  boxrule=0.6pt, arc=1pt, left=6pt, right=6pt
}
\newtcolorbox{pointcle}{
  enhanced, breakable,
  colback=blue!4, colframe=blue!55!black,
  fonttitle=\bfseries, title={Point clé},
  boxrule=0.8pt, arc=2pt, left=6pt, right=6pt
}
\newtcolorbox{piege}{
  enhanced, breakable,
  colback=red!4, colframe=red!60!black,
  fonttitle=\bfseries, title={Piège / difficulté},
  boxrule=0.8pt, arc=2pt, left=6pt, right=6pt
}
\newtcolorbox{commentaire}{
  enhanced, breakable,
  colback=gray!8, colframe=gray!55!black,
  fonttitle=\bfseries, title={Commentaire},
  boxrule=0.6pt, arc=2pt, left=6pt, right=6pt
}
\newtcolorbox{methode}{
  enhanced, breakable,
  colback=green!4, colframe=green!40!black,
  fonttitle=\bfseries, title={Méthode},
  boxrule=0.8pt, arc=2pt, left=6pt, right=6pt
}

\theoremstyle{definition}
\newtheorem*{solution}{Solution}

\title{\textbf{4MA305 -- Bases d'analyse fonctionnelle}\\[2mm]
\Large Révision des chapitres 1 à 3 (énoncés)}
\author{}
\date{Année universitaire 2026--2027}

\begin{document}
\maketitle

\begin{commentaire}
Feuille de révision des chapitres~1 à~3 du polycopié, dans l'optique de l'examen : deux exercices de questions de cours, deux vrai ou faux. L'idée est de chercher en conditions d'examen, sans document, avant de lire le corrigé détaillé (\texttt{corrige\_revision\_ch1-3.pdf}).
\end{commentaire}

\section*{Questions de cours}

% =====================================================================
\begin{exo}{Exercice 1 -- Questions de cours : chapitre 1}
\emph{À traiter sans document : énoncé complet, puis démonstration.}
\begin{enumerate}[label=(\alph*)]
  \item Énoncer et démontrer le théorème de prolongement des applications uniformément continues (Théorème~1.1).
  \item Énoncer et démontrer le lemme de Baire (Lemme~1.1).
  \item Énoncer et démontrer le théorème de Heine (Théorème~1.3).
  \item Montrer que dans un espace métrique complet, une partie est compacte si et seulement si elle est fermée et précompacte (Corollaire~1.4 ; on redémontrera les Propositions~1.13 et 1.14).
  \item Énoncer et démontrer le théorème de point fixe de Picard (Exercice~1.16).
\end{enumerate}
\end{exo}

% =====================================================================
\begin{exo}{Exercice 2 -- Questions de cours : chapitres 2 et 3}
\emph{À traiter sans document : énoncé complet, puis démonstration.}
\begin{enumerate}[label=(\alph*)]
  \item Énoncer et démontrer la caractérisation des applications linéaires continues entre EVN (Théorème~2.4).
  \item Énoncer et démontrer l'équivalence des normes sur $\K^d$ (Théorème~2.1).
  \item Montrer qu'un EVN est de Banach si et seulement si toute série absolument convergente y converge (Proposition~3.1 et Exercice~3.1).
  \item Énoncer et démontrer le théorème de Riesz (Théorème~3.1), en passant par le Lemme~3.1.
  \item Montrer que $\Lc(E,F)$ est un espace de Banach dès que $F$ en est un (Proposition~3.4).
  \item Énoncer et démontrer le théorème de Banach--Steinhaus (Théorème~3.7).
  \item Énoncer le théorème de l'application ouverte (Théorème~3.8, admis ici). En déduire le Corollaire~3.4 (continuité automatique de l'inverse), puis énoncer et démontrer le théorème du graphe fermé (Remarque~3.8).
\end{enumerate}
\end{exo}

\section*{Vrai ou faux}

% =====================================================================
\begin{exo}{Exercice 3 -- Vrai ou faux ? Chapitre 1}
Pour chaque affirmation, dire si elle est vraie ou fausse, et justifier par une preuve courte ou un contre-exemple. $(X,d)$ et $(Y,d')$ désignent des espaces métriques.
\begin{enumerate}[label=\arabic*.]
  \item L'adhérence de la boule ouverte $B(a,r)$ est la boule fermée $\Bf(a,r)$.
  \item Dans un espace métrique complet, toute partie fermée et bornée est compacte.
  \item L'image réciproque d'un compact par une application continue est compacte.
  \item L'image d'un fermé par une application continue est fermée.
  \item Une application continue transforme les suites de Cauchy en suites de Cauchy.
  \item Si $X$ est compact et $f:X\to Y$ est continue et bijective, alors $f^{-1}$ est continue.
  \item Si $X$ et $Y$ sont homéomorphes et $X$ est complet, alors $Y$ est complet.
  \item Dans un espace métrique complet, une intersection dénombrable d'ouverts denses est un ouvert dense.
  \item $\Q$ est une intersection dénombrable d'ouverts de $\R$.
  \item Tout espace métrique compact est séparable.
  \item Dans un espace métrique, toute partie fermée et précompacte est compacte.
  \item Si $D$ est dense dans $X$ et $f:D\to Y$ est uniformément continue, alors $f$ se prolonge en une application continue $X\to Y$.
\end{enumerate}
\end{exo}

% =====================================================================
\begin{exo}{Exercice 4 -- Vrai ou faux ? Chapitres 2 et 3}
Pour chaque affirmation, dire si elle est vraie ou fausse, et justifier par une preuve courte ou un contre-exemple. $E$ et $F$ désignent des EVN.
\begin{enumerate}[label=\arabic*.]
  \item Deux normes sur un même espace vectoriel sont toujours équivalentes.
  \item Si $E$ est de dimension finie, toute application linéaire $E\to F$ est continue.
  \item Tout sous-espace vectoriel de $E$ est fermé.
  \item Un sous-espace vectoriel strict de $E$ est d'intérieur vide.
  \item Si $T:E\to F$ est linéaire et bornée sur une boule $B(a,r)$, $r>0$, alors $T$ est continue.
  \item Pour $1\le p\le q\le\infty$, on a $L^p(]0,1[)\subset L^q(]0,1[)$.
  \item $\Cz([0,1])$ muni de la norme $\norm f_1=\int_0^1|f|$ est un espace de Banach.
  \item Un espace de Banach de dimension infinie n'admet pas de base (algébrique) dénombrable.
  \item La boule unité fermée de $\Cz([0,1])$ est compacte pour la norme uniforme.
  \item Une partie bornée et uniformément équicontinue de $\Cz([0,1])$ est compacte.
  \item Si $(f_n)$ est une suite de $\mathcal C^1([0,1])$ avec $\norm{f_n}_\infty+\norm{f_n'}_\infty\le M$ pour tout $n$, alors $(f_n)$ admet une sous-suite uniformément convergente.
  \item Si $E$ est de Banach et si $(T_n)$ est une suite de $\Lc(E,F)$ convergeant simplement vers $T$, alors $\opnorm{T_n-T}\to0$.
  \item Une limite simple d'applications linéaires continues de $E$ dans $F$ est continue.
  \item Une bijection linéaire continue $E\to F$ a un inverse continu.
\end{enumerate}
\end{exo}

\end{document}
